\section{Received Traditions Examined}

Great saints gather stories around them the way ancient cathedrals gather side
chapels: some rest on sixth-century foundations, while others express the
gratitude, theology, or iconography of much later ages. Six celebrated
traditions attached to Pope Saint Gregory the Great deserve plain, respectful
examination against the earliest controlled evidence.

\subsection{The Roman slave market: \emph{Non Angli sed Angeli}}

Few anecdotes in Christian history are better known than the story of Gregory
the monk pausing in the Roman marketplace before fair-skinned captive boys from
Britain, asking their nation, province, and king, and answering each name with
a Latin pun: Angles should be fellow-heirs with angels (\latin{Non Angli sed
Angeli}); men of Deira should be plucked from wrath (\latin{de ira}); and in
the kingdom of King \AE{}lle, \latin{Alleluia} must be sung---after which
Gregory allegedly set out in person for Britain until Pope Benedict I recalled
him at the clamor of the Roman crowd.

No hint of this scene appears anywhere in Gregory's surviving letters or
homilies. His earliest contemporary reference to the English, \work{Ep.}\ VI.7
(September 595), is a practical instruction to Candidus in Gaul to purchase
English slave youths of seventeen or eighteen for monastic formation; his 596
letters to the Frankish kings (\work{Ep.}\ VI.58--59) explain the mission
simply from reports that the English desired conversion while nearby clergy
neglected them. The slave-market encounter surfaces more than a century after
Gregory's death in the early eighth-century Whitby \work{Vita} and in Bede's
\work{Historia Ecclesiastica} II.1 (c.\ 731)---where Bede himself introduces it
cautiously as a tradition handed down from earlier generations---before John
the Deacon wove it into his ninth-century biography (reported in Barmby,
Prolegomena). As historical fact, the marketplace dialogue rests on early
received English tradition rather than contemporary testimony; yet
\work{Ep.}\ VI.7 proves that Gregory's real concern for English captives in
the slave trade was no invention, and the English story preserves with vivid
affection the memory of a pope who cared for their salvation before they ever
knew his name.

\subsection{The Archangel Michael atop Hadrian's Mausoleum}

During the plague of 590, medieval Roman legend relates that as Gregory led his
penitential procession toward Saint Peter's across the Tiber, he looked up to
the summit of the Emperor Hadrian's massive cylindrical tomb and beheld the
Archangel Michael sheathing a bloody sword in token that God's wrath had
ceased---giving the fortress its enduring name, Castel Sant'Angelo.

The contemporary eyewitness record refutes any early attestation of this
vision. Gregory of Tours (\work{Historia Francorum} X.1), writing directly from
his deacon who walked through Rome during that very procession, describes the
flood, the serpents, the exact seven starting churches, the convergence on
Saint Mary Major (not Saint Peter's), and eighty citizens falling dead in a
single hour---with complete silence regarding any angelic apparition. As Barmby
explicitly notes (Prolegomena), the vision is absent from all contemporaries
and from Bede as well. It took shape only in later medieval tradition
(popularized in thirteenth-century legendaries such as the \work{Golden
Legend}, linked to the early chapel of Saint Michael built atop the mausoleum).
It belongs to later devotional legend: unattested in 590, yet
remembered in Roman piety as a visual monument to Gregory's penitential
intercession in the city's darkest hour.

\subsection{The white dove whispering at Gregory's ear}

Across medieval sculpture, manuscript illumination, and stained glass, Saint
Gregory's unmistakable iconographic attribute is a white dove perched on his
shoulder, whispering into his ear as he dictates to his deacon Peter. The
narrative behind the image first appears in the eighth- and ninth-century
biographies (Paul the Deacon and John the Deacon), which tell how a secretary
peered behind a curtain during long pauses in Gregory's dictation of the
\work{Homilies on Ezekiel} and saw the Holy Spirit in the form of a dove
pressing its beak to the pope's lips.

Gregory's own prefaces offer a much more ordinary human picture: lecturing to
his monks at Constantinople, revising shorthand transcripts with scribes,
fighting bowel pain and slow fevers, and correcting pirated drafts
(\work{Moralia}, Ep.\ Ded.\ 2, 5; \work{Homiliae in Evangelia}, Praefatio).
The whispering dove is a later hagiographic and iconographic emblem,
not a contemporary biographical report. Read in its proper register, it states
in visual shorthand what the Church confessed of the Doctor's teaching: that
his scriptural exposition was guided by the Spirit of God.

\subsection{The ``Gregorian Mass'' (thirty consecutive Masses)}

Unlike the preceding legends, the Catholic devotional custom of the
``Gregorian Mass'' (or Trental)---having thirty consecutive daily Masses
celebrated for the repose of one deceased person---grows directly out of
Gregory's own pen, though with an important distinction between event and
institution. In \work{Dialogi} IV.55/57 (summarized in patristic scholarship
and cross-referenced in Migne PL 76), Gregory narrates that when a monk
named Justus at Saint Andrew's died concealing three gold coins against his
vow of poverty, Gregory imposed a stern burial outside the community grave to
warn the living, and then, thirty days later, moved with compassion, ordered
the prior Pretiosus to offer the Eucharistic sacrifice for Justus on thirty
consecutive days without interruption; at the end of the thirty days, Justus
appeared to his brother Copiosus announcing his release from suffering.

Gregory's \work{Homiliae in Evangelia} (37.7--8) confirms that this
deep faith in the liberating power of the Mass for the absent and departed was
thoroughly his own. Yet Gregory recorded one singular monastic act of mercy at
Saint Andrew's, not a universal liturgical law. The formalization of the
``Gregorian Masses'' into a standardized thirty-day devotional practice across
Latin Christendom developed gradually during the Middle Ages out of the
immense popularity of the \work{Dialogi}, building a medieval liturgical custom
upon an authentic Gregorian monastic narrative.

\subsection{Intercession for the pagan Emperor Trajan}

A dramatic early medieval legend relates that as Pope Gregory walked through
Trajan's Forum in Rome, he recalled the pagan Emperor Trajan's famous act of
justice---halting an imperial march to hear a poor widow's plea for her
murdered son---and wept so bitterly in Saint Peter's Basilica over the
emperor's damnation that God heard his prayer and delivered Trajan's soul.

Nothing of the kind stands in Gregory's authentic corpus. On the contrary, in
his \work{Homiliae in Evangelia} (40.7), commenting on the great
chasm between Lazarus and the rich man (Luke 16:26), Gregory states firmly
that ``neither do the unjust pass over to the lot of the blessed, because they
are bound by everlasting damnation, nor can the just pass over to the
reprobate, because, uplifted by the justice of judgment, they feel no
compassion to deliver them'' (\latin{Nec injusti ergo ad beatorum sortem
transeunt, quia damnatione perpetua constringuntur}). The Trajan story first
emerges in the early eighth-century Whitby \work{Vita} (ch.\ 29), passes
through Paul the Deacon and John the Deacon, and sparked centuries of subtle
scholastic debate---including Saint Thomas Aquinas's careful hypotheses in his
commentary on the \work{Sentences} and \work{De veritate} and Dante's
celebrated scenes in \work{Purgatorio} X and \work{Paradiso} XX. Historically,
the tale is an early medieval legend without sixth-century attestation;
theologically, it endured because medieval thinkers used Gregory's legendary
compassion to wrestle with the deep mystery of divine justice, invincible
ignorance, and intercessory prayer.

\subsection{Personal authorship of the Sacramentary and Gregorian chant}

As established in the preceding liturgical survey, the widespread medieval
belief that Gregory personally composed the entire Roman Sacramentary and the
full corpus of Latin plainchant rests on ninth-century attributions---most
influentially John the Deacon (\work{Vita} II.6, Migne PL 78, col.\ 649) and
the Carolingian rubrics of the \work{Hadrianum} (Wilson 1915). The internal
evidence of the \work{Hadrianum} itself---which contains Gregory's own feast
day marked off by the Carolingian \latin{Hucusque} preface---together with the
eighth- and ninth-century Frankish-Roman synthesis of plainchant shows that
these great monuments carry Gregory's name as patron, reformer, and anchor of
Roman liturgical identity rather than as their sole sixth-century author,
overlaying later attribution upon a genuine core of specific Gregorian reforms.
